\section{Introduction}\label{sec:introduction}
Which parameter bits should be protected when a neural network cannot be protected in full? Bit significance offers an inexpensive answer: suppress faults in the most significant positions first. Other methods use numerical perturbations, gradients, curvature, or measured fault effects to identify important locations. The engineering question is whether the extra information yields a better allocation under the fault probability and protection budget of interest.

The effect of a stored-bit fault depends on its encoding, location, input, and the other faults present. Consequently, a ranking that predicts isolated damage need not minimize loss when many bits can fail simultaneously. This observation does not make static protection obsolete. A fixed mask can be appropriate for a fixed operating condition, and a ranking can itself incorporate multi-fault information. What requires measurement is whether a particular allocation transfers across conditions.

We investigate this question with six frozen INT8 networks spanning vision, speech, anomaly detection, sentiment classification, and language modeling. We compare structural and gradient rankings throughout, and a published Hessian-based method where its computation has been validated. Each comparison uses the same proposed random fault fields, input cohort, eligible-bit population, and protection budget. Reliability is measured against logical bit-error probability (BER) and protected-bit fraction, not inferred from score correlations alone.

Three questions organize the evaluation: does a score predict isolated damage; does its resulting mask preserve task quality under simultaneous faults; and does the preferred allocation remain preferable when the budget or fault probability changes? These are different tests. A good answer to the first does not establish the second, and success at one operating point does not establish the third.

The contribution is a controlled comparison and its decision-level findings. Paired results reveal changes in relative method quality with model, BER, and budget. We also test whether conditioning on faulty backgrounds improves a local allocation beyond exact isolated-bit measurements. It often changes the mask, but a 32-pool KWS study does not establish a held-out improvement over that stronger comparator. This negative result limits the practical conclusion just as much as the observed ranking reversals.

Our earlier design flow included hardware synthesis and selective redundancy~\cite{guo2025inpixel}. The present study is software-only: it evaluates exact stored-code XOR faults and ideal suppression on protected coordinates. It neither measures physical mitigation overhead nor converts logical BER into a radiation rate. The intended outcome is an evidence-based choice of allocation for a declared operating condition, not another universally preferred score.
